\documentclass[12pt]{article}
\usepackage[utf8]{inputenc}
\usepackage[T1]{fontenc}
\usepackage{hyperref}
\title{Gambian Wolof Language Resources for K-12 Curriculum Development}
\author{Compiled by Claude Sonnet 5.5 (claude-sonnet-5-5) for Andr\'as Kornai}
\date{}
\begin{document}
\maketitle

\section{Introduction}
Gambian Wolof (ISO 639-3: wof; Glottocode: gamb1252) is an Atlantic-Congo language of wider communication in the Gambia. The reference description gives about 230,000 speakers and says it is not known to be taught in schools. One Gambian guide in the search results says Wolof and other national languages are taught in schools, which conflicts with that. The searches found mostly commercial dictionaries and learner books, not curriculum materials.

\section{Primary Educational Resources}
\begin{itemize}
\item \emph{Learn Wolof with Ease} (Balami, 2021) is aimed at newcomers to the Gambia. It covers everyday conversation, a concise grammar and a short lexicon. I saw no URL for it, so it is not in the JSON list.
\end{itemize}
No school textbooks or primers were found.

\section{Linguistic and Reference Materials}
\begin{itemize}
\item \emph{Wolof-English/English-Wolof Dictionary and Phrasebook} (Nyima Kantorek, Hippocrene). It is based on the Gambian dialect and has over 3,000 entries, a phrasebook and a pronunciation guide. \url{https://www.hippocrenebooks.com/store/p356/Wolof-English%2FEnglish-Wolof_Dictionary_%26_Phrasebook.html}
\item Cambridge Language Centre catalogue entry for the Kantorek dictionary: \url{https://www.langcen.cam.ac.uk/resources/langw/data/wo_107.html}
\item David P. Gamble's \emph{Gambian Wolof-English Dictionary} and \emph{Elementary Gambian Wolof Grammar} (both 1991, self-published) were described in the results. The author called the dictionary tentative. I saw no URL for either, so they are not in the JSON list.
\item Janga multilingual Wolof dictionaries (English-Wolof-French and English-French-Wolof) are sold through Polyglot Books. They are not specific to the Gambia. \url{https://payhip.com/PolyglotBooks/blog/articles/2-new-wolof-dictionaries-now-available}
\item A Wolof resources page for less commonly taught languages: \url{https://wisc.pb.unizin.org/lctlresources/?p=4079}
\end{itemize}

\section{Religious and Community Resources}
None found. The results mention Peace Corps volunteers compiling a Wolof and Mandinka dictionary and grammar, but I did not confirm they are available and saw no URL for them.

\section{Digital and Audio Resources}
\begin{itemize}
\item Gambia guide article on Wolof pronunciation: \url{https://www.my-gambia.com/mymagazine/learn-wollof-pronunciation/}
\item Gambia guide article on languages spoken in the Gambia: \url{https://www.my-gambia.com/mymagazine/languages-spoken-in-the-gambia/}
\end{itemize}
No online course or app specific to Gambian Wolof was found.

\section{Curriculum Development Context}
The existing material is mostly dictionaries and learner books for outsiders, and its spelling is partly custom (Kantorek). A curriculum would need an agreed orthography and new material written with Gambian educators. The question of whether Wolof is already taught in Gambian schools should be checked with local sources.

\end{document}